\chapter*{Remerciements}
\addcontentsline{toc}{chapter}{Remerciements}
\markboth{Remerciements}{}

Au terme de ce mini-projet, nous tenons à remercier toutes les personnes qui ont contribué à sa réalisation.

Nos premiers remerciements s'adressent à \textbf{Prof.~Maria Zemzami}, qui a proposé et encadré ce projet dans le cadre du module \enquote{Santé digitale et ingénierie des systèmes de santé}. Son cahier des charges~\cite{zemzami2026} (objectifs, livrables, normes à respecter et grille d'évaluation) a structuré notre démarche du début à la fin, et ce rapport en reprend volontairement le plan.

Nous remercions également le corps professoral de l'\ac{ENSAM} de Rabat, et particulièrement les enseignants de la filière \ac{INDIA}, pour la qualité de la formation dispensée~: les compétences en programmation, en modélisation et en sécurité des systèmes d'information acquises au cours de notre cursus ont été mobilisées à chaque étape du projet.

Enfin, chaque membre de l'équipe remercie les cinq autres pour la rigueur et l'entraide dont chacun a fait preuve~: une plateforme de téléconsultation n'est utile que si l'interface du patient, celle du médecin, le serveur, la vidéo et la sécurité parlent le même langage, et cet accord s'est construit ensemble.

\vspace{2cm}
\begin{flushright}
\textit{A.~\textsc{Issami}, M.~\textsc{Ezbakhe}, H.~\textsc{Benzahra},}\\
\textit{Z.~\textsc{Hsain}, H.~\textsc{Kichchou}, M.~\textsc{Hamdouch}}\\
\textit{2026}
\end{flushright}
